\chapter{Systembeschreibung und Durchführung}
\label{ch:durchfuehrung}

Dieses Kapitel beschreibt den entwickelten FTG-Stack im Detail,
den Integrationsprozess und die iterativen Schritte zur Fehlerdiagnose
und Kalibrierung. Wo Messdaten oder Fotos noch ausstehen,
sind entsprechende Platzhalter markiert.

% ============================================================
\section{Architektur des entwickelten FTG-Stacks}
% ============================================================

\subsection{Überblick der Pakete}

Der entwickelte Stack (\texttt{ftg\_mxck}) ergänzt den Basisstack
des MXCarkit um sieben spezialisierte ROS~2-Pakete. Der vollständige
Datenpfad ist in Abbildung~\ref{fig:pipeline} schematisch dargestellt.

\bildplatzhalter{Blockdiagramm: /scan $\to$ scan\_preprocessor $\to$
obstacle\_substitution $\to$ follow\_the\_gap\_v0 $\to$
ctu\_ftg\_adapter $\to$ ftg\_command\_node $\to$
/autonomous/ackermann\_cmd}{Verarbeitungspipeline des FTG-Stacks
(Blockdiagramm)}
\label{fig:pipeline}

\begin{table}[H]
\centering
\caption{ROS~2-Pakete des FTG-Stacks und ihre Rollen}
\label{tab:pakete}
\begin{tabularx}{\textwidth}{lX}
\toprule
\textbf{Paket} & \textbf{Aufgabe} \\
\midrule
\texttt{obstacle\_msgs}       & Benutzerdefinierte Nachrichtentypen für Hindernisdarstellung \\
\texttt{obstacle\_substitution} & Umwandlung jedes LiDAR-Strahls in ein Kreishindernis \\
\texttt{follow\_the\_gap\_v0} & CTU-FTG-Planerkern (C++): berechnet sicheren Durchfahrtswinkel \\
\texttt{mxck\_ftg\_perception} & Optionale LiDAR-Vorverarbeitung auf Frontbereich \\
\texttt{mxck\_ftg\_planner}   & Adapter: Winkel begrenzen, Geschwindigkeit nach Abstand wählen \\
\texttt{mxck\_ftg\_control}   & Steuerknoten: Ackermann-Kommandos aus Winkel + Geschwindigkeit \\
\texttt{mxck\_ftg\_bringup}   & Launch-Dateien zur koordinierten Inbetriebnahme \\
\bottomrule
\end{tabularx}
\end{table}

\subsection{Topic-Schnittstellen}

\begin{table}[H]
\centering
\caption{Kritische ROS~2-Topics in der Verarbeitungspipeline}
\label{tab:topics}
\begin{tabularx}{\textwidth}{llX}
\toprule
\textbf{Topic} & \textbf{Typ} & \textbf{Bedeutung} \\
\midrule
\texttt{/scan}                         & LaserScan       & Rohdaten des RPLIDAR \\
\texttt{/autonomous/ftg/scan\_filtered} & LaserScan       & Gefilterter Frontscan \\
\texttt{/obstacles}                    & ObstaclesStamped & Kreishindernisse \\
\texttt{/final\_heading\_angle}         & Float32         & FTG-Zielwinkel (rad) \\
\texttt{/gap\_found}                   & Bool            & Planerstatus \\
\texttt{/autonomous/ftg/gap\_angle}     & Float32         & Begrenzter Lenkwinkel (rad) \\
\texttt{/autonomous/ftg/target\_speed}  & Float32         & Zielgeschwindigkeit (m/s) \\
\texttt{/autonomous/ackermann\_cmd}     & AckermannDriveStamped & Finales Fahrkommando \\
\bottomrule
\end{tabularx}
\end{table}

\subsection{Konfigurationsparameter}

\begin{table}[H]
\centering
\caption{Wichtige Konfigurationsparameter des Stacks}
\label{tab:params}
\begin{tabularx}{\textwidth}{llX}
\toprule
\textbf{Parameter} & \textbf{Endwert} & \textbf{Beschreibung} \\
\midrule
\texttt{front\_center\_deg}        & $0°$      & LiDAR-Montagewinkel zur Fahrzeugfront \\
\texttt{front\_fov\_deg}           & $60°$     & Breite des ausgewerteten Frontbereichs \\
\texttt{clip\_min\_range\_m}       & $0{,}18$\,m & Minimale Messreichweite \\
\texttt{clip\_max\_range\_m}       & $2{,}0$\,m  & Maximale Messreichweite \\
\texttt{stop\_distance\_m}         & $0{,}30$\,m & Vollbremsung bei Unterschreitung \\
\texttt{slow\_distance\_m}         & $0{,}55$\,m & Langsamfahrt bei Unterschreitung \\
\texttt{clear\_distance\_m}        & $0{,}90$\,m & Freie Fahrt bei Überschreitung \\
\texttt{cruise\_speed\_mps}        & $0{,}25$\,m/s & Reisegeschwindigkeit \\
\texttt{steering\_limit\_deg}      & $22°$     & Maximaler Lenkwinkel \\
\texttt{gap\_angle\_limit\_rad}    & $0{,}60$\,rad & Maximaler Lückenwinkel \\
\texttt{enable\_turn\_speed\_scaling} & false  & Drehwinkelbasierte Geschwindigkeitsreduktion \\
\bottomrule
\end{tabularx}
\end{table}

% ============================================================
\section{Identifizierte Ausgangsprobleme und deren Behebung}
% ============================================================

Bei der Analyse des bestehenden Stacks vor der eigentlichen Integration
wurden mehrere architekturelle Probleme identifiziert:

\subsection{Problem~P1: Potenziell doppelter Planer-Output}

Die ursprüngliche Launch-Datei \texttt{ftg\_stack.launch.py} erlaubte
es, sowohl den CTU-Adapterpfad als auch einen alternativen Planerpfad
(\texttt{ftg\_planner\_node}) gleichzeitig zu starten. Bei gleichzeitigem
Betrieb beider Pfade würden mehrere Knoten auf denselben Output-Topics
publizieren, was zu undefiniertem Steuerverhalten führt.

\textbf{Lösung:} Einführung des Launch-Parameters \texttt{use\_ftg\_planner}
(Standard: \texttt{false}) mit gegenseitigem Ausschluss -- entweder der
CTU-Adapter \emph{oder} der alternative Planner-Knoten wird gestartet.

\subsection{Problem~P2: Gestapelte Geschwindigkeitspolitik}

Die Standardkonfiguration aktivierte gleichzeitig die clearance-basierte
Geschwindigkeitswahl im \texttt{ctu\_ftg\_adapter\_node}
\emph{und} eine zusätzliche lenkwinkelbasierte Geschwindigkeitsskalierung
im \texttt{ftg\_command\_node}. Das führte dazu, dass die Geschwindigkeit
an zwei verschiedenen Stellen reduziert wurde, was die Kalibrierung
deutlich erschwerte.

\textbf{Lösung:} Deaktivierung der redundanten Skalierung durch Setzen von
\texttt{enable\_turn\_speed\_scaling: false} in \texttt{ftg\_control.yaml}.

% ============================================================
\section{Integration und Kompilierung}
\label{sec:integration}
% ============================================================

Der Stack wurde im Docker-Container \texttt{mxck2\_development}
(ROS~2 Foxy, Python~3.8, CMake) kompiliert. Der Kompilierungsbefehl:

\begin{lstlisting}[language=bash, caption={Build-Befehl für den FTG-Stack}]
colcon build --symlink-install \
  --packages-select \
    obstacle_msgs \
    obstacle_substitution \
    follow_the_gap_v0 \
    mxck_ftg_perception \
    mxck_ftg_planner \
    mxck_ftg_control \
    mxck_ftg_bringup
\end{lstlisting}

Die Kompilierung dauerte auf dem Jetson Nano ca.\ 3\,min\,50\,s.
Das Paket \texttt{obstacle\_msgs} benötigte die längste Kompilierzeit
($\approx$\,2\,min\,12\,s) aufgrund der ROS~2-Nachrichtengenerierung.
Es traten \emph{Clock-Skew}-Warnungen auf, die auf eine
Zeitdiskrepanz zwischen Host-Dateisystem und Container zurückzuführen
sind, die Funktionalität jedoch nicht beeinträchtigten.

Das System wird gestartet mit:

\begin{lstlisting}[language=bash, caption={Systemstart via Launch-Datei}]
ros2 launch mxck_ftg_bringup ftg_full_system.launch.py \
  use_vehicle_control:=false \
  record_bag:=false
\end{lstlisting}

Beim Start werden sechs ROS~2-Prozesse initiiert:
\texttt{robot\_state\_publisher}, \texttt{scan\_preprocessor\_node},
\texttt{obstacle\_substitution\_node}, \texttt{follow\_the\_gap},
\texttt{ctu\_ftg\_adapter\_node} und \texttt{ftg\_command\_node}.

% ============================================================
\section{Iterativer Kalibrierungsprozess}
\label{sec:kalibrierung}
% ============================================================

Die korrekte Kalibrierung des Parameters \texttt{front\_center\_deg}
erwies sich als zentrales Problem der Integration. Dieser Parameter
gibt an, bei welchem Winkel im Rohlaserscan die Fahrzeugvorderseite liegt.
Im Folgenden werden die durchgeführten Optimierungsiterationen beschrieben.

\subsection{Iteration~0: Initialer Start -- keine Lücken gefunden}

Beim ersten Start des Systems lieferte der FTG-Planerkern
kontinuierlich folgende Meldung:

\begin{lstlisting}[caption={Ausgabe des FTG-Kerns beim ersten Start}]
[follow_the_gap-4] No gaps found
[follow_the_gap-4] No gaps found
...
\end{lstlisting}

Der \texttt{/autonomous/ftg/scan\_filtered}-Topic zeigte, dass
der gesamte Frontbereich mit dem Minimalwert $0{,}05$\,m belegt war,
also keinerlei gültige Messdaten enthielt.

\textbf{Diagnose:} Der Kalibrierungswert \texttt{front\_center\_deg = 135°}
entsprach nicht der tatsächlichen Montagelage des LiDAR-Sensors.
Der ausgewertete ``Frontbereich'' lag physikalisch \emph{hinter}
oder seitlich des Fahrzeugs. Zusätzlich fehlte die TF-Transformation
\texttt{base\_link} $\to$ \texttt{laser}, weshalb der Preprocessor
keinen automatischen Offset berechnen konnte.

\subsection{Iteration~1: front\_center\_deg = 90°}

\textbf{Änderung:} \texttt{front\_center\_deg} von $135°$ auf $90°$ geändert.

\textbf{Ergebnis:} Der Scan-Preprocessor lieferte nun Daten,
jedoch zeigte \texttt{/final\_heading\_angle} konstant $\approx -1{,}83$\,rad
-- unabhängig von der Hindernisposition:

\begin{table}[H]
\centering
\caption{Messwerte bei front\_center\_deg = 90° (Iteration~1)}
\label{tab:iter1}
\begin{tabular}{lll}
\toprule
\textbf{Hindernisposition} & \textbf{/final\_heading (rad)} & \textbf{/gap\_angle (rad)} \\
\midrule
Vorne mittig & $-1{,}811$ & $-0{,}600$ (Grenzwert) \\
Vorne rechts & $-1{,}809$ & $-0{,}600$ (Grenzwert) \\
Vorne links  & $-1{,}809$ & $-0{,}600$ (Grenzwert) \\
\bottomrule
\end{tabular}
\end{table}

Der Wert $-0{,}600$ entspricht exakt dem Parameter
\texttt{gap\_angle\_limit\_rad}: Der Adapter begrenzte den Winkel stets
auf diesen Maximalwert, weil der Planerkern permanent $-1{,}83$\,rad ausgab.

\textbf{Diagnose:} Der Frontbereich lag noch immer in der falschen Richtung.
Die Obstacle-Koordinaten für ein Hindernis direkt vor dem Fahrzeug
(ca.\ 35\,cm Distanz) zeigten:

\begin{lstlisting}[caption={Obstacle-Koordinaten bei Hindernis vorne (Iteration~1)}]
center: x=0.336, y=0.862  (Hindernis "vorne mittig")
center: x=0.345, y=0.872  (Hindernis "vorne rechts")
center: x=0.328, y=0.904  (Hindernis "vorne links")
\end{lstlisting}

Alle drei Positionen liegen in der gleichen $x/y$-Region,
das System unterscheidet nicht zwischen links, mitte und rechts.
Der hohe $y$-Wert zeigt, dass die Fahrzeugvorderseite im
Koordinatensystem entlang der $y$-Achse statt der $x$-Achse zeigt --
ein klassischer Achsendrehfehler.

\subsection{Iteration~2: Verengung des Sichtfelds}

\textbf{Änderungen:}
\begin{itemize}
    \item \texttt{front\_fov\_deg}: $100°$ $\to$ $60°$
    \item \texttt{clip\_max\_range\_m}: $5{,}0$\,m $\to$ $2{,}0$\,m
    \item \texttt{clip\_min\_range\_m}: auf $0{,}18$\,m (echtes LiDAR-Minimum)
\end{itemize}

\textbf{Beobachtung:} Das Fahrzeug zeigte erstmals Bewegungen:
Es lenkte abwechselnd links und rechts und fuhr manchmal kurz vor,
wenn viel freier Raum vorhanden war. Dies entsprach jedoch
\emph{nicht} dem erwarteten Ausweichverhalten.

\bildplatzhalter{Videostandbild oder Foto aus dem Experiment:
Fahrzeug lenkt in die falsche Richtung anstatt auszuweichen}{
Beobachtetes Fehlverhalten: Fahrzeug reagiert in entgegengesetzte Richtung}

\subsection{Iteration~3: Vollständige Kalibrierung (front\_center\_deg = 0°)}

\textbf{Änderung:} \texttt{front\_center\_deg} von $90°$ auf $0°$ geändert.

\textbf{Ergebnis:} Das System lieferte nun plausible, positionsabhängige
Lenkwinkel:

\begin{table}[H]
\centering
\caption{Messwerte nach vollständiger Kalibrierung (Iteration~3)}
\label{tab:iter3}
\begin{tabular}{llll}
\toprule
\textbf{Hindernis} & \textbf{/final\_heading (rad)} & \textbf{/gap\_angle (rad)} & \textbf{Clearance (m)} \\
\midrule
Karton vorne ($\approx 60$\,cm)   & $+0{,}179$ & $+0{,}182$ & $0{,}906$ \\
Wand $2$\,m vorne                & $+0{,}179$ & $+0{,}179$ & $0{,}910$ \\
Hindernis vorne rechts           & $+0{,}181$ & $+0{,}178$ & $0{,}910$ \\
Hindernis vorne links            & $+0{,}180$ & $+0{,}258$ & --        \\
\bottomrule
\end{tabular}
\end{table}

\datenplatzhalter{Foxglove-Screenshot oder Plot von /final\_heading\_angle
und /gap\_angle über Zeit, während ein Hindernis von links nach rechts
bewegt wird -- zeigt ob die Winkelreaktion korrekt ist.}

Der gefilterte Scan zeigte einen sauberen, engbegrenzten Frontbereich:
\begin{lstlisting}[caption={Scan-Metadaten nach korrekter Kalibrierung}]
angle_min: -0.522 rad  (ca. -30 Grad)
angle_max: +0.522 rad  (ca. +30 Grad)
range_min: 0.18 m
range_max: 2.0 m
\end{lstlisting}

% ============================================================
\section{Praktischer Fahrttest}
% ============================================================

\subsection{Manueller Fahrttest}

Zur Verifikation der Aktuatoranbindung wurde das Fahrzeug zunächst
manuell per RC-Fernbedienung gesteuert. Die Motorgeschwindigkeit
(\texttt{/commands/motor/speed}) zeigte dabei Werte zwischen
$700$--$950$\,RPM, was dem normalen Betriebsbereich des VESC entspricht.

Die Ausgabe auf \texttt{/autonomous/ackermann\_cmd} beim manuellen Fahren:
\begin{lstlisting}[caption={Ackermann-Kommando während manueller Fahrt}]
drive: { steering_angle: 0.0, speed: 0.1,
         acceleration: 0.0, jerk: 0.0 }
\end{lstlisting}

\subsection{Autonomer Fahrttest}

Im autonomen Modus reagierte das Fahrzeug auf Hindernisse.
Bei einem Frontobjekt in $\approx 0{,}41$\,m Abstand:
\begin{lstlisting}[caption={Planerstatus bei nahem Frontobjekt}]
gap_found=True, heading_fresh=True, scan_fresh=True,
gap_angle_rad=+0.4788, front_clearance_m=0.408,
target_speed=0.180 m/s, reason=crawl_distance
\end{lstlisting}

Das Fahrzeug fuhr kurz geradeaus mit reduzierter Geschwindigkeit
(Kriechmodus: $0{,}18$\,m/s) und lenkte leicht zur Seite.
Es stieß jedoch gegen das Hindernis, bevor eine vollständige
Kurskorrektur stattfinden konnte.

\bildplatzhalter{Videostandbilder oder Bildsequenz aus dem
autonomen Fahrttest: (1) Hindernis erkannt, (2) Lenkreaktion,
(3) Kollision}{Ablauf des autonomen Fahrttests}

\subsection{Aufgezeichnetes Rosbag}

Zur späteren Analyse wurde ein 15-sekündiges Rosbag aufgezeichnet,
in dem ein Hindernis vor dem Fahrzeug von links nach rechts bewegt wurde:

\begin{lstlisting}[language=bash, caption={Rosbag-Aufnahmebefehl}]
ros2 bag record -s mcap -o /mxck2_ws/bags/ftg_diag_small \
  /scan /autonomous/ftg/scan_filtered /obstacles \
  /final_heading_angle /gap_found \
  /autonomous/ftg/gap_angle \
  /autonomous/ftg/planner_status /tf_static
\end{lstlisting}

\datenplatzhalter{Foxglove-Wiedergabe des aufgezeichneten Rosbags:
Plot von /final\_heading\_angle und /gap\_found über Zeit +
3D-Ansicht der /obstacles-Punkte. Zeigt die Reaktion des Systems
auf die Hindernisbewegung.}
